% 第六章 链上音乐资源发现支撑模块设计与实现

\chapter{链上音乐资源发现与检索机制设计}


\section{模块目标与总体设计}

在 FreeFlow 系统中，资源发现模块承担链上音乐资源检索、定位和结果组织等支撑任务。链上音乐资源同时具有传统音乐信息与链上身份信息两类属性。前者包括标题、艺术家、专辑、流派和描述等元数据，后者包括链网络编号、合约地址、Token ID、发布交易哈希、IPFS 内容标识以及资源唯一键等标识字段。两类属性在检索场景中的作用存在明显差异：传统元数据主要服务于歌曲发现与内容浏览，链上身份字段主要服务于精确定位、资产确认和链上资源追踪。基于这一特点，系统中的资源发现模块需要同时支持自然语言查询与链上身份查询，并在统一结果列表中给出清晰的结果依据。

本模块面向系统冷启动阶段的资源发现需求，主要关注以下几个方面。当用户输入歌曲名、艺术家、专辑名或流派标签时，系统能够召回语义相关的链上音乐资源；当用户输入 Token ID、合约地址、交易哈希、IPFS CID 或 resourceKey 等链上标识时，系统能够优先返回目标明确的资源；当查询中存在大小写差异、空格变化、符号干扰或轻微拼写偏差时，系统仍能保持较好的召回能力与结果稳定性；同时，系统还需要在结果展示时保留字段命中情况和得分来源，以支持结果分析、功能调试与用户侧解释。

结合上述目标，本文将链上音乐资源发现支撑模块划分为查询归一化、候选召回、规则评分和结果排序四个阶段。该结构能够将查询处理、候选生成与结果组织解耦，便于后续根据系统规模和数据积累情况继续扩展更复杂的检索能力。


\section{检索对象与索引字段设计}

\subsection{检索对象定义}

系统以已发布的\texttt{CreatorRelease}作为检索对象。考虑到链上音乐资源需要具备可确认的链上身份，系统仅将状态为\texttt{PUBLISHED}且\texttt{chainId}、\texttt{musicAssetAddress}和\texttt{tokenId}均非空的资源纳入检索范围。

该约束保证了检索结果具备明确的链上可验证性，能够与后续播放、购买、确权和展示等业务环节形成一致的数据基础。

\subsection{索引字段划分}

为兼顾自然语言检索与链上身份检索，本文将检索字段划分为文本字段、身份字段和辅助排序字段三类。各字段用途如表\ref{tab:ranking_fields_rewrite}所示。

\begin{table}[htbp]
    \linespread{1.5}
    \zihao{5}
    \songti
    \centering
    \caption{链上音乐资源检索字段设计}\label{tab:ranking_fields_rewrite}
    \begin{tabular}{|p{3.35cm}|p{8.6cm}|}
        \hline
        字段                                                       & 用途                    \\ \hline
        \texttt{title}                                           & 标题匹配、标题词元重合计算与模糊相似度计算 \\ \hline
        \texttt{artistName}                                      & 艺术家名称匹配               \\ \hline
        \texttt{albumName}                                       & 专辑名称匹配                \\ \hline
        \texttt{genreLabel}                                      & 流派标签匹配                \\ \hline
        \texttt{description}                                     & 长文本辅助召回               \\ \hline
        \texttt{chainId}                                         & 链网络识别                 \\ \hline
        \texttt{musicAssetAddress}                               & 音乐资源合约定位              \\ \hline
        \texttt{platformHubAddress}                              & 平台相关合约定位              \\ \hline
        \texttt{tokenId}                                         & 链上作品唯一 Token 定位       \\ \hline
        \texttt{publishTxHash}                                   & 发布交易定位                \\ \hline
        \texttt{audioCid}、\texttt{coverCid}、\texttt{metadataCid} & IPFS 内容寻址检索           \\ \hline
        \texttt{resourceKey}                                     & 资源唯一键定位               \\ \hline
        \texttt{comments}、\texttt{purchases}                     & 行为热度估计                \\ \hline
        \texttt{publishedAt}、\texttt{updatedAt}                  & 时效性估计                 \\ \hline
    \end{tabular}
\end{table}

上述字段设计体现出两项原则。其一，标题、艺术家、专辑与流派承担主要的内容发现功能；其二，Token ID、合约地址、CID 和交易哈希承担高确定性的资源定位功能。通过对这两类字段的联合组织，系统能够在歌曲搜索与链上资产搜索两类场景下保持较稳定的表现。


\section{查询处理与候选召回}

\subsection{查询归一化}

由于用户输入可能受到大小写、分隔符、拷贝格式以及轻微拼写偏差的影响，系统在召回前首先对原始查询进行归一化处理。归一化流程包括以下几个步骤：

\begin{enumerate}
    \item 对原始查询执行 Unicode NFKC 归一化，以统一全角半角字符和兼容字符表示；
    \item 将字母统一为小写形式；
    \item 合并多余空格并去除首尾空白；
    \item 按非字母数字分隔符切分词元；
    \item 生成用于紧凑比较的\texttt{compact}文本，以支持跨空格和符号差异的匹配。
\end{enumerate}

在归一化基础上，系统还会识别几类具有显著格式特征的链上输入，包括 EVM 地址、IPFS CID、Token ID、交易哈希和\texttt{resourceKey}。当查询被识别为此类字段时，后续结果组织阶段会提高身份字段的匹配优先级，以增强精确定位能力。

本文当前实现采用字符串归一化与轻量词元切分方案，尚未引入完整中文分词与向量语义检索模型。结合原型系统当前的资源规模和冷启动背景，这一实现方式能够覆盖大部分典型查询场景，并保持较好的可复现性与解释性，便于后续分析与调整。

\subsection{候选召回策略}

候选召回采用“字段召回 + 近期补充”的两阶段策略。设用户查询为$q$，候选集合记为$C$。系统首先基于归一化后的查询执行数据库字段召回，获得候选集合$C_1$；随后引入一批近期发布资源作为补充集合$C_2$，并与$C_1$合并去重，得到最终候选集合$C$。

其中，$C_1$主要通过数据库中的\texttt{contains}条件构造，召回目标字段包括标题、艺术家、专辑、流派、描述以及链上身份字段。该阶段能够较高效地覆盖精确匹配、包含匹配和部分格式变化场景。

考虑到仅依赖 SQL 字段包含匹配时，轻微拼写误差和非标准输入可能导致漏召回，系统进一步引入近期发布资源作为补充集合$C_2$。这一处理能够在候选范围内保留一定数量的新资源，为后续模糊比较提供可参考对象。对于冷启动平台而言，近期资源本身也具有较高的业务关注度，因此该补充方式能够在召回鲁棒性与系统资源开销之间取得较好的平衡。

为控制后续处理阶段的计算复杂度，系统对候选规模设置上限。当前实现中，主召回候选数根据查询命中情况动态调整，总量限制在 500 条以内；补充召回的近期资源不超过 200 条。该限制能够保证资源发现链路在常规负载下保持可控延迟。


\section{结果组织规则与评分机制}

\subsection{评分机制定义}

为便于对候选资源进行统一比较，系统采用固定权重的规则化加和方式生成结果分值，并将该实现记为\texttt{chain\_resource\_rank\_v1}。对任意候选资源$d$和查询$q$，其总分定义如公式\eqref{eq:rank_score_rewrite}所示：

\begin{equation}
    Score(d,q)=S_{identity}(d,q)+S_{text}(d,q)+S_{fuzzy}(d,q)+S_{freshness}(d)+S_{popularity}(d)
    \label{eq:rank_score_rewrite}
\end{equation}

其中，$S_{identity}$表示身份匹配分，$S_{text}$表示文本匹配分，$S_{fuzzy}$表示模糊相似度分，$S_{freshness}$表示时效性分，$S_{popularity}$表示热度分。

该评分结构主要服务于结果组织过程，可对应为三类判断信息。第一类为资源定位信息，用于判断候选资源与查询在链上身份层面的对应关系；第二类为内容相关信息，用于判断候选资源与查询在标题、艺术家和专辑等内容层面的接近程度；第三类为辅助区分信息，用于在多个相关候选之间进一步比较发布时间和用户关注程度。

\subsection{身份匹配规则}

身份匹配规则面向链上资源的高确定性定位需求，覆盖\texttt{resourceKey}、\texttt{tokenId}、\texttt{musicAssetAddress}、\texttt{publishTxHash}以及各类 CID 字段。此类字段具有唯一性强、歧义度低、定位能力高的特点，因此在总分中占据较高比重。

当前实现中，\texttt{resourceKey}、\texttt{tokenId}、音乐合约地址和 IPFS CID 的精确匹配权重分别设置为 100、80、70 和 65。其权重设置依据主要来自字段的业务作用与定位强度：\texttt{resourceKey}直接对应系统内部资源标识，具有最高区分度；\texttt{tokenId}在指定合约上下文中具有明确唯一性；合约地址能够稳定定位到资源集合或资产发行主体；CID 能够定位到链下内容对象，在链上音乐场景中同样具有较强的检索价值。

\subsection{文本匹配规则}

文本匹配规则覆盖标题、艺术家、专辑、流派和描述五类字段。系统分别考虑精确匹配、前缀匹配、包含匹配和词元重合四种信号，并通过加权方式体现不同字段与不同匹配关系的重要程度。

当前实现中的文本字段权重如表\ref{tab:text_matching_weights}所示。

\begin{table}[htbp]
    \linespread{1.5}
    \zihao{5}
    \songti
    \centering
    \caption{文本字段匹配权重设置}\label{tab:text_matching_weights}
    \begin{tabular}{|l|c|c|c|c|}
        \hline
        字段  & 精确匹配 & 前缀匹配 & 包含匹配 & 词元重合 \\ \hline
        标题  & 45   & 30   & 20   & 14   \\ \hline
        艺术家 & 34   & 24   & 16   & 12   \\ \hline
        专辑  & 24   & 16   & 10   & 8    \\ \hline
        流派  & 20   & 14   & 8    & 6    \\ \hline
        描述  & 10   & 8    & 6    & 4    \\ \hline
    \end{tabular}
\end{table}

上述设置体现了查询意图在不同字段上的重要性差异。标题通常最直接反映用户的目标作品，因此具有较高权重；艺术家和专辑承担较强的检索引导作用；流派更多服务于主题浏览和风格发现；描述字段文本长度较长、表述方式不稳定，其作用主要体现在辅助召回和弱相关区分。

\subsection{模糊匹配规则}

为增强系统对轻微拼写错误、分隔符变化和局部字符缺失的适应能力，结果组织阶段进一步引入模糊相似度规则。本文采用基于 Levenshtein 距离的归一化字符串相似度，对查询与候选资源的\texttt{compact}文本进行比较。设归一化后的查询紧凑文本为$q_c$，候选字段紧凑文本为$d_c$，则相似度可表示为公式\eqref{eq:lev_similarity}：

\begin{equation}
    sim_{lev}(q_c,d_c)=1-\frac{Lev(q_c,d_c)}{\max(|q_c|,|d_c|)}
    \label{eq:lev_similarity}
\end{equation}

其中，$Lev(\cdot)$表示编辑距离，$|\cdot|$表示字符串长度。相似度越高，表示二者在字符层面越接近。

为降低短字符串导致的误召回风险，系统仅在\texttt{compact}文本长度不小于 3 时启用模糊相似度评分；当查询长度较短时，系统采用更高的相似度阈值。该约束能够抑制由极短词、缩写和单字符差异造成的不稳定排序现象。

\subsection{新鲜度与热度规则}

当多个候选资源都满足文本相关或身份相关条件时，系统进一步引入时效性与热度信息进行区分。

新鲜度分采用指数衰减形式。设资源发布时间距当前时刻的天数为$ageDays$，则新鲜度定义如公式\eqref{eq:freshness_rewrite}所示：

\begin{equation}
    freshness(d)=0.5^{ageDays/45}
    \label{eq:freshness_rewrite}
\end{equation}

该式表示资源每经过 45 天，其新鲜度贡献衰减为原值的一半。当前实现中，新鲜度项的最高贡献为 8 分。该设计能够使近期发布资源在分值接近时获得更高排序位置，同时避免时间因素过度主导结果顺序。

热度分由评论数与购买数共同决定，其归一化形式如公式\eqref{eq:popularity_rewrite}所示：

\begin{equation}
    popularity(d)=\min\left(1,\frac{\ln(1+comments+3\times purchases)}{\ln(31)}\right)
    \label{eq:popularity_rewrite}
\end{equation}

当前实现中，热度项的最高贡献为 7 分。式中对购买数给予更高权重，原因在于购买行为代表了更强的资源确认信号与价值投入，相较评论行为具有更高的可信度和业务指向性。对数归一化处理能够减弱极端高活跃样本对结果的影响，使热度特征在不同规模资源之间保持更平滑的区分能力。


\section{结果说明与处理流程}

\subsection{结果说明机制}

为了增强检索结果的可分析性，系统在输出结果时同步记录各项得分来源。对于每个候选资源，系统会保留命中的字段类型、匹配关系、各项特征分值及其对总分的贡献，以及新鲜度与热度的归一化得分。其中，字段类型可以是标题、艺术家、Token ID 或 CID 等；匹配关系可以是精确匹配、前缀匹配、包含匹配、词元重合或模糊相似。通过保留上述信息，系统能够在结果明细中展示排序依据，从而支持误召回分析、参数调整和结果解释。

例如，当某一资源因标题精确匹配获得 45 分、艺术家前缀匹配获得 24 分、新鲜度获得 6.3 分、热度获得 4.8 分时，系统能够在结果明细中展示上述组成，使结果来源具备可追踪性。该机制有助于分析误召回、定位参数设置问题，并为后续功能完善提供参考。

\subsection{处理流程}

链上音乐资源发现流程如图\ref{fig:search_ranking_pipeline}所示。系统先对查询执行归一化与特殊字段识别，再通过数据库召回与近期补充获得候选集合，随后依次计算身份匹配、文本匹配、模糊相似度、新鲜度和热度信息，最终根据总分完成结果排序并输出解释信息。

\begin{figure}[H]
    \centering
    \includegraphics[width=\textwidth]{figures/out/search_ranking_pipeline.png}
    \caption{链上音乐资源发现处理流程图}
    \label{fig:search_ranking_pipeline}
\end{figure}

对应的处理过程如算法\ref{algo:chain_resource_rank_rewrite}所示。

\begin{algorithm}[H]
    \caption{链上音乐资源发现处理流程}
    \label{algo:chain_resource_rank_rewrite}
    \small
    \begin{algorithmic}[1]
    \Require 用户查询$q$，返回数量$k$
    \Ensure 排序后的链上音乐资源列表
    \State 对$q$进行归一化处理，得到词元集合、\texttt{compact}文本和特殊字段识别结果
    \State 基于标题、艺术家、专辑、流派、描述及链上身份字段执行数据库召回，得到候选集合$C_1$
    \State 选取近期发布资源构成补充集合$C_2$
    \State 合并去重，得到最终候选集合$C=C_1\cup C_2$
    \For{每个候选资源$d \in C$}
        \State 分别计算身份匹配分、文本匹配分、模糊相似度分、新鲜度分和热度分
        \State 汇总总分$Score(d,q)$并记录结果解释信息
    \EndFor
    \State 过滤未命中有效相关特征的候选资源
    \State 按总分降序排序；若总分相同，则按发布时间降序排序
    \State 返回前$k$个结果
    \end{algorithmic}
\end{algorithm}

\section{复杂度分析}

设候选资源数量为$N$，平均参与比较的文本长度为$L$。候选召回阶段主要由数据库字段查询完成，计算开销主要集中在结果组织阶段。

文本匹配分的计算主要依赖字段比较与词元重合统计，其时间复杂度近似为$O(NL)$。模糊相似度采用 Levenshtein 距离，在单次比较中的时间复杂度近似为$O(L^2)$，因此在候选集合上的总体复杂度近似为$O(NL^2)$。最终排序采用基于总分的比较排序，其复杂度为$O(N\log N)$。

考虑到编辑距离计算对长文本较为敏感，系统在实现中对参与模糊比较的文本长度进行截断，并对候选数量设置上限，因此整体运行开销能够维持在可接受范围内。对于原型系统当前的资源规模，该复杂度设计能够满足离线评估与在线检索的基本需求。


\section{检索效果评估设计}

为验证资源发现模块在冷启动阶段的支撑效果，本文构建了一个离线评估集。该评估集包含 56 条模拟链上音乐资源和 62 个查询，按照 14 个艺术家或专辑簇组织，用于覆盖多种具有代表性的检索情形。

在查询类型方面，评估集包括中英文标题检索、长专辑名检索、流派检索，以及 Token ID、合约地址、IPFS CID、\texttt{resourceKey}和交易哈希等链上身份检索。该设计能够同时考察系统对自然语言查询和链上标识查询的处理能力。

在相关性标注方面，单目标定位类查询采用唯一三级相关资源标注方式，适用于标题、Token ID、CID 和\texttt{resourceKey}等高确定性场景；艺术家、专辑和流派等多结果查询采用三级相关性标注，按照查询意图匹配程度和资源代表性标记为 3、2、1 三个等级。该标注方式能够更细致地反映多结果场景下的结果优劣。

当前评估集属于面向原型验证的结构化模拟数据集，主要用于检验冷启动阶段资源发现链路的工程可用性与结果合理性。基于真实线上行为日志构建的大规模评估体系仍需在系统投入实际运行后进一步完善。

本文从排序质量、召回有效性和响应效率三个方面评估检索效果。其中，nDCG@10 的使用参考了分级相关性条件下的经典累积增益评价方法\cite{Jarvelin2002CG}。其中，Precision@10 用于衡量前 10 个结果中的相关资源比例，MRR 用于衡量首个高相关结果出现的位置，nDCG@10 用于衡量分级相关性条件下前 10 个结果的排序质量，零结果率用于反映检索请求中无返回结果的比例，平均响应时间则用于衡量检索链路的效率表现。

对比方法设置为四组：数据库字段包含匹配基线、仅使用文本评分的方法、文本评分叠加链上身份字段的方法，以及完整的\texttt{chain\_resource\_rank\_v1}实现。通过分层对比，可以观察身份字段、模糊相似度以及辅助排序信息对整体检索效果的影响。具体实验过程、数据分布和结果分析见第 7 章表\ref{tab:ranking_query_categories}、表\ref{tab:ranking_metrics_actual}和图\ref{fig:ranking_latency_curve}。


\section{本章小结}

本章围绕链上音乐资源发现支撑模块的设计与实现展开论述，完成了检索对象定义、索引字段划分、查询归一化、候选召回、规则评分、结果解释机制和评估方案设计。该模块面向系统冷启动阶段的资源发现需求，能够在缺乏大规模行为数据的条件下支持文本检索、链上身份定位和结果说明等功能，为后续章节中的实验验证提供实现基础。
